\documentclass[11pt]{article}
\usepackage[T1]{fontenc}
\usepackage{lmodern,geometry,xcolor,hyperref}
\geometry{letterpaper,margin=1in}
\definecolor{navy}{HTML}{243746}
\hypersetup{pdftitle={Timed narration: recommendations in matching markets},
pdfauthor={Ishan Khillan, Elaine Zhang, Todd Zhou}}
\setlength{\parindent}{0pt}
\setlength{\parskip}{9pt}
\renewcommand{\familydefault}{\sfdefault}
\begin{document}

{\small Harvard CS 2880 / AI for Social Impact}\par
\section*{Slide 1: Who gets a place?}
{\color{navy}\textbf{00:00--00:18 | 18 seconds}}\par\medskip
We study how AI recommendations change access to opportunities with limited capacity. A useful list for one person can change someone else's outcome. Our question is simple: when everyone follows the recommendations, who gets a place?
\par
\newpage
{\small Harvard CS 2880 / AI for Social Impact}\par
\section*{Slide 2: A better suggestion can hurt someone else}
{\color{navy}\textbf{00:18--01:11 | 53 seconds}}\par\medskip
Start with two applicants and two school seats. Applicant A prefers school x to school y. Applicant B accepts only x. School x gives A priority over B.
\par
Initially, A is shown y and B is shown x. The blue lines show their assignments. Both get a seat.
\par
Now change only A's recommendation. A sees x instead of y. Both applicants still see exactly one school. But they now compete for x, and A has higher priority. A gets x, B gets no seat, and y stays empty.
\par
The new suggestion helps A. B's own list has not changed at all, but B loses access. That is the basic reason to evaluate recommendations jointly. A person's outcome depends on what other people are encouraged to pursue. This can happen in housing and jobs as well as schools.
\par
\newpage
{\small Harvard CS 2880 / AI for Social Impact}\par
\section*{Slide 3: Choose a shortlist, then assign the places}
{\color{navy}\textbf{01:11--02:06 | 55 seconds}}\par\medskip
We separate the recommendations from the assignment rule. First, the recommender chooses a shortlist of feasible options for each person. Then the assignment rule decides who gets which place.
\par
We measure three things. A match means someone receives a place. Target coverage means a particular reference choice appears on their shortlist. A lost match means someone receives a place under one policy, but not under the policy we compare it with.
\par
The reference differs across our studies. In rentals it is the simulated assignment when all feasible options are available. In Chile it is the school recorded in later enrollment.
\par
Our question is whether recommendations can improve while everyone who had a place stays matched somewhere. They may switch schools. Keeping a seat does not necessarily mean getting a better school.
\par
\newpage
{\small Harvard CS 2880 / AI for Social Impact}\par
\section*{Slide 4: Same personal shortlists. Twice the match rate.}
{\color{navy}\textbf{02:06--02:56 | 50 seconds}}\par\medskip
Here is our clearest finding. In the rental simulation, focus on renters with only one to five feasible apartments. Every policy already shows them all their options.
\par
Their match rate rises from 13.1 percent to 27.5 percent when we change how recommendations are coordinated across the market. That is more than double. Yet their own shortlists stay exactly the same. Only other renters' shortlists change.
\par
We compare a policy coordinating one option with a policy balancing the whole shortlist. I will explain both shortly. The result shows how much competition can matter even when your own recommendations never change. It is a simulation result, not measured housing outcomes.
\par
\newpage
{\small Harvard CS 2880 / AI for Social Impact}\par
\section*{Slide 5: Why study rentals and school admissions?}
{\color{navy}\textbf{02:56--04:07 | 71 seconds}}\par\medskip
The settings serve different purposes.
\par
Rentals let us study actual saved GPT recommendations and compare ways to coordinate them. Listings and household profiles come from market and survey data, but renter preferences, landlord priorities, and resulting assignments are modeled.
\par
We reconstruct 1,280 markets and verify all 55,115 cached candidate hashes. These new comparisons use previously analyzed calendar windows. Some homes recur across them, so this is a follow-up analysis, not independent property-level validation.
\par
Chile gives us families' submitted rankings and the actual admissions rules. Our reconstruction matches the released course or no-offer outcomes in about two million audited applicant-years. That checks institutional fidelity.
\par
The new shortlists still produce counterfactual outcomes under fixed reports. The connection across settings is competition for limited capacity, not an assumption that housing and schools behave identically.
\par
\newpage
{\small Harvard CS 2880 / AI for Social Impact}\par
\section*{Slide 6: Two ways to coordinate recommendations}
{\color{navy}\textbf{04:07--05:18 | 71 seconds}}\par\medskip
We test two simple coordination approaches.
\par
The anchor first computes a maximum matching on all feasible pairs. It prefers pairs already on the shortlists when several maximum matchings are available. Each assigned option is then inserted into its recipient's shortlist, replacing at most one option. List lengths stay fixed.
\par
Balancing takes a broader approach. Count how many times each listing is recommended. Minimize the sum of those counts squared, which discourages everyone from seeing the same few listings. Among equally balanced shortlists, favor options higher in the saved GPT ranking. We solve this exactly with min-cost flow.
\par
Both policies use feasible pairs and saved recommendations. Neither sees the realized preference shocks, landlord priorities, or reference assignments.
\par
Balancing can replace several options per person. It is a stronger coordination control, not an equal-edit-budget comparison with the one-option anchor. Neither approach imposes an individual retention constraint.
\par
\newpage
{\small Harvard CS 2880 / AI for Social Impact}\par
\section*{Slide 7: Individual coverage misses the placement loss}
{\color{navy}\textbf{05:18--06:13 | 55 seconds}}\par\medskip
First look at the original comparison with random shortlists of the same length. Blue is reference-option coverage. Orange is placement. Zero means no difference from random.
\par
Across all three enriched calendar windows, GPT increases coverage by about eight percentage points. Placement falls by about two to four points. The directions repeat, and the original market-bootstrap intervals exclude zero.
\par
This does not mean GPT makes worse predictions about everything. It means that covering an individual's full-list assignment does not determine how many people can be placed when everyone's shortlists interact.
\par
These markets deliberately include affordable listings and renters with feasible options. The ordinary sampling design and all uncertainty details are in the appendix.
\par
\newpage
{\small Harvard CS 2880 / AI for Social Impact}\par
\section*{Slide 8: Coordination changes the comparison}
{\color{navy}\textbf{06:13--07:20 | 67 seconds}}\par\medskip
Now put four policies on the same axes. Farther right means more reference coverage. Higher means more people placed.
\par
Random shortlists are on the left. GPT moves right, but slightly down. The one-option anchor moves GPT upward, recovering placement with a small change in coverage.
\par
Balanced shortlists go much higher. Relative to the anchor, they gain another 12.6 percentage points of placement while losing about three points of reference coverage in this panel.
\par
The comparison matters. Relative to plain random, balancing improves both outcomes. So the evidence does not show an unavoidable conflict between coverage and placement.
\par
These are point means for the September and October enriched panel. A control using extra information from the simulation improves both measures further. Detailed controls and paired intervals are in the appendix. We have not established GPT as the best recommender. This plot also cannot tell us whether any individual is protected.
\par
\newpage
{\small Harvard CS 2880 / AI for Social Impact}\par
\section*{Slide 9: A net gain does not tell us who gained}
{\color{navy}\textbf{07:20--08:14 | 54 seconds}}\par\medskip
Return to the anchor and compare it with the original GPT shortlists on exactly the same markets.
\par
There are 449 newly placed renter appearances. There are also 37 exits: people who were placed before and are now unplaced. Subtracting the exits from the entrants gives the 412 net gain.
\par
Reporting only that net gain would hide the loss of access. The more effective balancing policy has the same issue. Compared with the anchor, it creates many more placements but also causes 171 exits.
\par
An exit is defined relative to a named policy. These counts describe repeated simulated appearances, not unique households. Their role is to test an individual protection claim that an average placement rate cannot establish.
\par
\newpage
{\small Harvard CS 2880 / AI for Social Impact}\par
\section*{Slide 10: Why can an unchanged shortlist help you?}
{\color{navy}\textbf{08:14--09:10 | 56 seconds}}\par\medskip
The cleanest competition result concerns renters with only one to five feasible listings. Our policies show up to five options, so these renters already see everything. Their personal shortlists cannot change.
\par
When the whole market uses the anchor, this group's placement rate is 13.1 percent. Under balancing, it is 27.5 percent. That is a 14.4-point gain even though the group's own recommendations are identical.
\par
The change comes from other renters' shortlists. Coordinating those shortlists changes who competes for each listing.
\par
This is a sizable positive spillover inside the model. It still does not protect every member of the group: 15 previously placed appearances lose access under balancing. Better group outcomes and zero individual losses remain different claims.
\par
\newpage
{\small Harvard CS 2880 / AI for Social Impact}\par
\section*{Slide 11: Chile: the losses are concentrated}
{\color{navy}\textbf{09:10--10:42 | 92 seconds}}\par\medskip
Rentals give us controlled comparisons, but their preferences are modeled. Chile lets us examine the same access question using families' submitted rankings and a reconstruction of the admissions rules.
\par
We compare two classifiers trained on an earlier cohort to predict initial offers. One uses submitted rank, and the other does not. They select short shortlists for later cohorts. We score coverage against later observed enrollment and replay admissions on the selected options.
\par
The rank-informed policy gains 4,638 enrollment hits overall, but loses 1,670 placements relative to rank-blind shortlists.
\par
The graph separates pupils by whether they historically received an initial offer. Each bar uses its own group's denominator. Pupils with no initial offer make up 7.4 percent of the linked sample but account for 76 percent of the net placement loss.
\par
That is 1,270 net lost placements in the smaller group. Its placement rate falls by 0.873 percentage points in the counterfactual comparison.
\par
These are not revoked official offers. Historical offer status defines the groups; placement loss compares two new shortlist policies under fixed reported preferences. The concentration of the loss is what motivates the protection question.
\par
\newpage
{\small Harvard CS 2880 / AI for Social Impact}\par
\section*{Slide 12: The comparison policy matters}
{\color{navy}\textbf{10:42--11:37 | 55 seconds}}\par\medskip
There is an important check on that result. Compare the same rank-informed policy with each family's first five discretionary courses instead.
\par
Now it improves both aggregates: 9,800 more enrollment hits and 8,143 more placements. Again, a prediction--placement conflict is not inevitable.
\par
But the group without a historical initial offer still loses 2,860 counterfactual placements against this alternative.
\par
That is why the comparison policy and the affected population both need to be visible. An overall improvement can coexist with worse access for a particular group. We should not treat the historical enrollment target as welfare, or assume that a larger total resolves the distributional question.
\par
\newpage
{\small Harvard CS 2880 / AI for Social Impact}\par
\section*{Slide 13: Choosing shortlists with zero exits is hard}
{\color{navy}\textbf{11:37--12:38 | 61 seconds}}\par\medskip
Here is the mathematical design problem.
\par
Fix the baseline shortlists and a supplied set of priority draws. Choose one new shortlist policy that works across every draw. Maximize extra reference hits, subject to every baseline recipient keeping a seat on every draw. Recipients may switch schools, but cannot become unplaced.
\par
In Chile, we start from rank-blind shortlists, preserve options shared with rank-informed shortlists, fill from their union, and keep each quota.
\par
Our theorem says that deciding whether a target gain is achievable is NP-complete, even with unit seats and shortlists of length at most three. Checking a chosen policy with deferred acceptance remains polynomial.
\par
The theorem uses fixed, strict priorities on an explicit panel. Chile's family-linked rules are tested separately. The difficulty is choosing compatible edits, as the next slide illustrates.
\par
\newpage
{\small Harvard CS 2880 / AI for Social Impact}\par
\section*{Slide 14: Why individually safe edits can conflict}
{\color{navy}\textbf{12:38--13:34 | 56 seconds}}\par\medskip
The proof starts with a graph. Each vertex represents a beneficial shortlist edit. Each edge creates one priority draw where the two endpoint edits can conflict.
\par
Look at this one draw. Selecting edit u displaces another, already-matched pupil toward the fallback seat. Selecting edit v displaces a different already-matched pupil toward the same seat. Either edit alone is safe. Select both, and one pupil must lose a seat. Other vertex edits are harmless on this draw.
\par
We create a draw for every edge. One fixed set of edits must work on all of them. So a safe set cannot contain both endpoints of any edge. That is an independent set.
\par
Each selected vertex adds one reference hit. Maximizing safe gains therefore solves the Independent Set problem. The appendix gives the full construction.
\par
\newpage
{\small Harvard CS 2880 / AI for Social Impact}\par
\section*{Slide 15: Allowing school switches leaves more room}
{\color{navy}\textbf{13:34--14:42 | 68 seconds}}\par\medskip
We then compare two protection requirements on the same Chilean design panels: 1,024 lotteries per cohort for 2022 through 2024.
\par
If every baseline recipient must keep the same course and seat type, the maximum additional enrollment coverage is exactly 314 hits within our allowed shortlist family.
\par
If recipients may switch schools while keeping a seat, we construct shortlists gaining 642 hits. The certified upper bound is 4,620. Allowing switches therefore permits more than twice the demonstrated coverage gain.
\par
But the bracket is wide. We have not found the optimal retention policy.
\par
These are hindsight optimization benchmarks, using the evaluation cohorts' later enrollment labels. They do not describe a policy that could choose its shortlists before those outcomes existed. To make the constructive answer useful, we also need a rule based only on information available before assignment.
\par
\newpage
{\small Harvard CS 2880 / AI for Social Impact}\par
\section*{Slide 16: A test panel is weaker than a guarantee}
{\color{navy}\textbf{14:42--15:33 | 51 seconds}}\par\medskip
We evaluate two rules frozen after applications arrive, before the ordinary assignment. Neither uses that cohort's realized offers or later enrollment to choose shortlists.
\par
The structural rule protects a set of interactions around applicants without a guaranteed fallback. Its proof covers all permitted ordinary lottery orders under the fixed inputs and reconstructed mechanism.
\par
It gains 28 enrollment hits and has zero exits in the test panels.
\par
A learned rule passes its design-panel checks and gains 50 hits. But it causes an exit in each of two fresh lottery generators.
\par
The guaranteed gain is small. The useful distinction is between a proof covering the allowed orders and a policy that merely passed many simulated draws.
\par
\newpage
{\small Harvard CS 2880 / AI for Social Impact}\par
\section*{Slide 17: Protecting recipients can preserve exclusion}
{\color{navy}\textbf{15:33--16:36 | 63 seconds}}\par\medskip
The social objective cannot stop with protecting people who already have a place. Our structural Chile rule produces no new entrants on the reported paired test lotteries. It preserves incumbents, but does not expand access for previously unplaced pupils in those tests.
\par
That raises a substantive design choice. How should we balance protection for existing recipients against opportunities for people who start excluded?
\par
Our measurements are incomplete answers. Historical enrollment reflects constraints as well as preference. Receiving any seat can mean moving to a worse school, and modeled ranks are not measured household welfare.
\par
Before deployment, people should retain the ability to add choices, correct errors, and seek human review or appeal. Those are proposed safeguards, not interventions we have tested. Recommendations redistribute existing opportunities. They cannot create affordable apartments or new school capacity.
\par
\newpage
{\small Harvard CS 2880 / AI for Social Impact}\par
\section*{Slide 18: Evaluate the allocation, not just the list}
{\color{navy}\textbf{16:36--17:06 | 30 seconds}}\par\medskip
Recommendations need a market-level evaluation.
\par
Coordination can improve both coverage and placement. But better totals do not guarantee that each person keeps access.
\par
We measure that distinction, formalize the retention problem, and test the gap between a finite test panel and a guarantee.
\par
The next design problem is to improve access for people without a place while protecting those who already have one. Thank you.
\par
\end{document}
